% GENERATED FILE. Edit canonical lessons, then run npm run generate:books.
% canonical-source-hash: fnv1a64:b19a557b55c30710
% canonical-lessons: PA-C187-past-return, PA-C187-si, PA-C187-si-review, PA-C187-san, PA-C187-si-san-contrast, PA-C187-san-description, PA-C187-san-habitual, PA-C187-san-review, PA-C187-payoff

\chapter{He Was There; They Were There}
\label{ch:pa-a2-past-copula}
\hlchaptermodality{187}

\begin{chapteropening}
\textbf{By the end of this chapter:} \emph{I can recognise and use common singular and plural past-copula forms in descriptions and familiar routines.}

Retrieve exact singular and plural past-copula forms through listening, speaking, reading, and bounded copying after a slow contrast ramp.
\end{chapteropening}

\section[(past copula return)]{Start with one verb you know}

\label{lesson:PA-C187-past-return}



\begin{quote}
Say the Punjabi infinitive for “to be.”
\end{quote}

\begin{grammarlens}[title={Grammar Lens: one past contrast at a time}]
Keep the familiar meaning steady. The next step will make the singular past form. After a short retrieval, a separate step will add the plural and polite past form.
\end{grammarlens}

\subsection*{Your turn}
\begin{itemize}
\raggedright
  \item \emph{Hear:} \emph{honaa}; YOU SAY: \textbf{\pa{ਹੋਣਾ}}
\end{itemize}

\begin{tcolorbox}[breakable,colback=teal!4,colframe=teal!35!black,title={Before you move on}]
Say “to be” once more in Punjabi.

This voice-first review introduces no new written target.
\end{tcolorbox}
\section[\texorpdfstring{\pa{ਸੀ} (sii)}{ਸੀ (sii)}]{\pa{ਸੀ} — was}

\label{lesson:PA-C187-si}



\begin{quote}
Say \textbf{\pa{ਹੋਣਾ}}, then \textbf{\pa{ਅਮਨ}}, \textbf{\pa{ਘਰ}}, and \textbf{\pa{ਵਿੱਚ}}.
\end{quote}

\begin{grammarlens}[title={Grammar Lens: the singular past form}]
\begin{quote}
\textbf{\pa{ਸੀ}} — \emph{sii} — was
\end{quote}

Use this past copula with a singular subject here. Keep the short form whole; the plural and polite past form comes after one retrieval step.
\end{grammarlens}

\subsection*{Your turn}
\begin{itemize}
\raggedright
  \item \emph{Hear:} \emph{sii}; YOU SAY: \textbf{\pa{ਸੀ}}
  \item \emph{Read it:} \textbf{\pa{ਅਮਨ} \pa{ਘਰ} \pa{ਵਿੱਚ} \pa{ਸੀ}}
  \item \emph{Write it:} copy \textbf{\pa{ਸੀ}} once and check the final \textbf{\pa{ੀ}}
\end{itemize}

\begin{tcolorbox}[breakable,colback=teal!4,colframe=teal!35!black,title={Before you move on}]
Read the exact singular past form. (\textbf{\pa{ਸੀ}}.)

Source: standard Punjabi uses \textbf{\pa{ਸੀ}} as a singular past copula.
\end{tcolorbox}
\section[\texorpdfstring{\pa{ਅਮਨ} \pa{ਘਰ} \pa{ਵਿੱਚ} \pa{ਸੀ} (Aman ghar vich sii)}{ਅਮਨ ਘਰ ਵਿੱਚ ਸੀ (Aman ghar vich sii)}]{Hold the singular past form steady}

\label{lesson:PA-C187-si-review}



\begin{quote}
Say the exact form for “was” once without looking.
\end{quote}

\subsection*{Your turn}
\begin{itemize}
\raggedright
  \item \emph{Hear:} \emph{Aman ghar vich sii}; YOU SAY: \textbf{\pa{ਅਮਨ} \pa{ਘਰ} \pa{ਵਿੱਚ} \pa{ਸੀ}}
  \item \emph{Read it:} \textbf{\pa{ਅਮਨ} \pa{ਘਰ} \pa{ਵਿੱਚ} \pa{ਸੀ}}
  \item \emph{Write it:} cover the line and write \textbf{\pa{ਸੀ}}
\end{itemize}

The meaning is past and the subject is singular. No second past form is needed yet.

\begin{tcolorbox}[breakable,colback=teal!4,colframe=teal!35!black,title={Before you move on}]
Retrieve \textbf{\pa{ਸੀ}} one more time.

This review adds no new past form.
\end{tcolorbox}
\section[\texorpdfstring{\pa{ਸਨ} (san)}{ਸਨ (san)}]{\pa{ਸਨ} — were}

\label{lesson:PA-C187-san}



\begin{quote}
Say \textbf{\pa{ਸੀ}}, then \textbf{\pa{ਫੁੱਲ}} and \textbf{\pa{ਲਾਲ}}.
\end{quote}

\begin{grammarlens}[title={Grammar Lens: the plural and polite past form}]
\begin{quote}
\textbf{\pa{ਸਨ}} — \emph{san} — were
\end{quote}

Use this past copula with a plural subject here. It also serves respectful singular reference, matching Punjabi’s familiar plural-for-respect pattern.
\end{grammarlens}

\subsection*{Your turn}
\begin{itemize}
\raggedright
  \item \emph{Hear:} \emph{san}; YOU SAY: \textbf{\pa{ਸਨ}}
  \item \emph{Read it:} \textbf{\pa{ਫੁੱਲ} \pa{ਲਾਲ} \pa{ਸਨ}}
  \item \emph{Write it:} copy \textbf{\pa{ਸਨ}} once and check the final \textbf{\pa{ਨ}}
\end{itemize}

\begin{tcolorbox}[breakable,colback=teal!4,colframe=teal!35!black,title={Before you move on}]
Say “was,” then “were.” (\textbf{\pa{ਸੀ}, \pa{ਸਨ}}.)

Source: standard Punjabi uses \textbf{\pa{ਸਨ}} as the plural or polite past copula.
\end{tcolorbox}
\section[\texorpdfstring{\pa{ਸੀ} · \pa{ਸਨ} (sii / san)}{ਸੀ · ਸਨ (sii / san)}]{One was; more than one were}

\label{lesson:PA-C187-si-san-contrast}



\begin{quote}
Say the pair once without looking.
\end{quote}

\subsection*{Your turn}
\begin{itemize}
\raggedright
  \item \emph{Hear:} \emph{Aman ghar vich sii}; YOU SAY: \textbf{\pa{ਅਮਨ} \pa{ਘਰ} \pa{ਵਿੱਚ} \pa{ਸੀ}}
  \item \emph{Hear:} \emph{phull laal san}; YOU SAY: \textbf{\pa{ਫੁੱਲ} \pa{ਲਾਲ} \pa{ਸਨ}}
  \item \emph{Read it:} \textbf{\pa{ਸੀ} · \pa{ਸਨ}}
  \item \emph{Write it:} cover the pair and write \textbf{\pa{ਸੀ}, \pa{ਸਨ}}
\end{itemize}

Keep the time past in both lines. Change the copula when the subject changes from singular to plural.

\begin{tcolorbox}[breakable,colback=teal!4,colframe=teal!35!black,title={Before you move on}]
Which form is singular? Which is plural or polite?

This review adds no perfective verb form.
\end{tcolorbox}
\section[\texorpdfstring{\pa{ਫੁੱਲ} \pa{ਲਾਲ} \pa{ਸਨ} (phull laal san)}{ਫੁੱਲ ਲਾਲ ਸਨ (phull laal san)}]{Keep “were” inside a small description}

\label{lesson:PA-C187-san-description}



\begin{quote}
Say the plural past form once.
\end{quote}

\subsection*{Your turn}
\begin{itemize}
\raggedright
  \item \emph{Hear:} \emph{phull laal san}; YOU SAY: \textbf{\pa{ਫੁੱਲ} \pa{ਲਾਲ} \pa{ਸਨ}}
  \item \emph{Read it:} \textbf{\pa{ਫੁੱਲ} \pa{ਲਾਲ} \pa{ਸਨ}}
  \item \emph{Write it:} copy \textbf{\pa{ਸੀ}, \pa{ਸਨ}} and underline \textbf{\pa{ਸਨ}}
\end{itemize}

The plural subject \textbf{\pa{ਫੁੱਲ}} keeps the plural past copula \textbf{\pa{ਸਨ}}.

\begin{tcolorbox}[breakable,colback=teal!4,colframe=teal!35!black,title={Before you move on}]
Say the singular and plural pair again.

This description reuses familiar words and adds no new colour or noun.
\end{tcolorbox}
\section[\texorpdfstring{\pa{ਉਹ} \pa{ਕੰਮ} \pa{ਕਰਦੇ} \pa{ਸਨ} (oh kamm karde san)}{ਉਹ ਕੰਮ ਕਰਦੇ ਸਨ (oh kamm karde san)}]{Put \pa{ਸਨ} after a familiar plural habit}

\label{lesson:PA-C187-san-habitual}



\begin{quote}
Say \textbf{\pa{ਉਹ} \pa{ਕੰਮ} \pa{ਕਰਦੇ}}, then retrieve \textbf{\pa{ਸਨ}}.
\end{quote}

\begin{grammarlens}[title={Grammar Lens: a familiar habit moved into the past}]
\begin{quote}
\textbf{\pa{ਉਹ} \pa{ਕੰਮ} \pa{ਕਰਦੇ} \pa{ਸਨ}} — \emph{oh kamm karde san} — they used to work
\end{quote}

The familiar plural habitual predicate stays before \textbf{\pa{ਸਨ}}. The final copula moves the routine into the past.
\end{grammarlens}

\subsection*{Your turn}
\begin{itemize}
\raggedright
  \item \emph{Hear:} \emph{oh kamm karde san}; YOU SAY: \textbf{\pa{ਉਹ} \pa{ਕੰਮ} \pa{ਕਰਦੇ} \pa{ਸਨ}}
  \item \emph{Read it:} \textbf{\pa{ਉਹ} \pa{ਕੰਮ} \pa{ਕਰਦੇ} \pa{ਸਨ}}
  \item \emph{Write it:} copy \textbf{\pa{ਕਰਦੇ} \pa{ਸਨ}} once
\end{itemize}

\begin{tcolorbox}[breakable,colback=teal!4,colframe=teal!35!black,title={Before you move on}]
Complete \textbf{\pa{ਉਹ} \pa{ਕੰਮ} \pa{ਕਰਦੇ} …} with the past copula.

This step reuses the already-owned habitual predicate and introduces no new perfective form.
\end{tcolorbox}
\section[\texorpdfstring{\pa{ਸੀ} · \pa{ਸਨ} (sii / san)}{ਸੀ · ਸਨ (sii / san)}]{Two past forms, three familiar frames}

\label{lesson:PA-C187-san-review}



\begin{quote}
Say “was; were” in that order.
\end{quote}

\subsection*{Your turn}
\begin{itemize}
\raggedright
  \item \emph{Hear:} \emph{Aman ghar vich sii}; YOU SAY: \textbf{\pa{ਅਮਨ} \pa{ਘਰ} \pa{ਵਿੱਚ} \pa{ਸੀ}}
  \item \emph{Hear:} \emph{phull laal san}; YOU SAY: \textbf{\pa{ਫੁੱਲ} \pa{ਲਾਲ} \pa{ਸਨ}}
  \item \emph{Hear:} \emph{oh kamm karde san}; YOU SAY: \textbf{\pa{ਉਹ} \pa{ਕੰਮ} \pa{ਕਰਦੇ} \pa{ਸਨ}}
  \item \emph{Write it:} cover the lines and write \textbf{\pa{ਸੀ}, \pa{ਸਨ}}
\end{itemize}

One singular description takes \textbf{\pa{ਸੀ}}. Both plural frames take \textbf{\pa{ਸਨ}}.

\begin{tcolorbox}[breakable,colback=teal!4,colframe=teal!35!black,title={Before you move on}]
Retrieve both exact forms without looking.

This review adds no perfective verb and no future form.
\end{tcolorbox}
\section[\texorpdfstring{\pa{ਅਮਨ} \pa{ਘਰ} \pa{ਵਿੱਚ} \pa{ਸੀ} … (Aman ghar vich …)}{ਅਮਨ ਘਰ ਵਿੱਚ ਸੀ … (Aman ghar vich …)}]{A gentle past-copula payoff}

\label{lesson:PA-C187-payoff}



\begin{quote}
Without looking, say the singular past form and then the plural or polite past form.
\end{quote}

\subsection*{Your turn: four skills, bounded to the two forms}
\begin{itemize}
\raggedright
  \item \emph{Hear:} \emph{Aman ghar vich sii; phull laal san; oh kamm karde san}; YOU SAY: \textbf{\pa{ਅਮਨ} \pa{ਘਰ} \pa{ਵਿੱਚ} \pa{ਸੀ}; \pa{ਫੁੱਲ} \pa{ਲਾਲ} \pa{ਸਨ}; \pa{ਉਹ} \pa{ਕੰਮ} \pa{ਕਰਦੇ} \pa{ਸਨ}}
  \item \emph{Read it:} \textbf{\pa{ਅਮਨ} \pa{ਘਰ} \pa{ਵਿੱਚ} \pa{ਸੀ।} \pa{ਫੁੱਲ} \pa{ਲਾਲ} \pa{ਸਨ।} \pa{ਉਹ} \pa{ਕੰਮ} \pa{ਕਰਦੇ} \pa{ਸਨ।}}
  \item \emph{Write it:} cover the model and write \textbf{\pa{ਸੀ}, \pa{ਸਨ}}
\end{itemize}

Underline the singular form once and the plural form twice.

\begin{tcolorbox}[breakable,colback=teal!4,colframe=teal!35!black,title={Before you move on}]
Retrieve both measured forms once more.

This is the tenth prerequisite tranche exposed by Punjabi A2 mock drafting. It does not unblock Punjabi A2 mock 1 by itself. It is not an A2 mock, not A2 readiness or attainment, and not external or human validation.

Source: the standard Punjabi past-copula contrast applied to familiar singular and plural phrases.
\end{tcolorbox}
